一维无限深势阱初始区间为 $0<x<L$，质量为 $m$ 的粒子最初处于基态 $\psi(x,0^-)=\sqrt{2/L}\sin(\pi x/L)$。在 $t=0$ 突然将右墙移至 $x=2L$，随后新势阱保持不变。移动足够快，可使用突然近似：$\psi(x,0^+)=\psi(x,0^-)$（并在 $L<x<2L$ 延拓为 $0$）。
(1) 求新阱的本征能量与初态在各新本征态上的展开系数 $c_n$；(2) 求测得新阱能级 $n=2$ 的概率（$n=2$ 是第一激发态）和所有偶数 $n\ne2$ 的概率；(3) 写出 $t>0$ 波函数并计算骤变后能量期望，解释与旧能量的关系；(4) 判断新哈密顿量能量方差是否有限，并说明原因。
